% LaTeX fragment.  Uses standard theorem/lemma/proof environments.
% Prepared independently, with the finite telescoper in the repository's
% Stirling draft checked and extended from integers to positive real names.
\section{Real convexity, effective uniqueness, and gamma inequalities}
\label{sec:real-convexity}

Throughout this section an argument $x>0$ is supplied with rational bounds
$0<a\le x\le b$.  Every real-valued expression denotes its fixed primitive
recursive rational approximation procedure.  An assertion with arbitrary
such procedures is a theorem schema, not quantification over an internally
available universe of real functions.  Integral notation below abbreviates
finite rational quadrature with the indicated error estimates.

\subsection{Digamma estimates and strong log-convexity}

Put $g(x)=\log\Gamma(x)$ and $\psi(x)=g'(x)$.  Differentiating the finite
Euler approximants gives
\[
 g_N'(x)=\log N-\sum_{k=0}^{N}\frac1{x+k},\qquad
 g_N''(x)=\sum_{k=0}^{N}\frac1{(x+k)^2}.
\]
The product construction supplies their uniform convergence on $[a,b]$.
For the second derivatives a sufficient tail estimate is
\[
 0\le\sum_{k=N+1}^{M}\frac1{(x+k)^2}
 \le\int_N^M\frac{dt}{(x+t)^2}\le\frac1{x+N}.
\]
Consequently
\begin{equation}\label{eq:psi-prime-positive}
 \psi'(x)=\sum_{k=0}^{\infty}\frac1{(x+k)^2},\qquad
 \frac1{b^2}\le\psi'(x)\le\frac1{a^2}+\frac1a
 \quad(a\le x\le b).
\end{equation}
To justify differentiation without a general interchange theorem, apply
the finite fundamental theorem to $g_N'$ on a rational segment, replace its
second derivative by a partial sum, and use the displayed uniform tail.
The resulting finite increment identities have an error tending to zero
at the displayed primitive recursive rate.  The same argument for $g_N$
identifies its first derivative.

For completeness, comparison of $t\mapsto(x+t)^{-1}$ with its upper and
lower rectangle sums gives
\[
 \log x-\frac1x\le\psi(x)\le\log x.
\]
The upper bound can be sharpened.  Convexity of this reciprocal function
gives, on each unit interval, that its integral is at most its trapezoidal
sum.  Summing and substituting in the formula for $g_N'$ yields
\begin{equation}\label{eq:digamma-elementary-bounds}
 \log x-\frac1x\le\psi(x)\le\log x-\frac1{2x}.
\end{equation}
All discarded terminal terms are bounded by rational multiples of $1/N$;
for example $0\le\log(1+x/N)\le b/N$.

Let $a\le x,y\le b$ and $0\le s\le1$.  The function
$g(t)-t^2/(2b^2)$ has nonnegative second derivative.  Integrating its
nondecreasing first derivative on the two sides of a division point proves
the secant inequality, giving
\begin{equation}\label{eq:strong-log-convexity}
 g((1-s)x+sy)
 \le(1-s)g(x)+sg(y)-\frac{s(1-s)(x-y)^2}{2b^2}.
\end{equation}
One can derive the same conclusion directly by adding finite interval
increments; no mean-value point is chosen.  Exponentiating proves
log-convexity of gamma.  The explicit last term proves strictness whenever
$s$, $1-s$, and $|x-y|$ have positive rational lower bounds.

In particular gamma is strictly convex as well: differentiating once more
gives $\Gamma''=\Gamma(\psi^2+\psi')>0$.  This conclusion has an explicit
positive gap on each interval $[a,b]$, using a positive lower bound for
$\Gamma$ there and \eqref{eq:psi-prime-positive}.

\subsection{A Bohr--Mollerup squeeze which needs no floor of a real}

Let $F$ be a fixed positive primitive recursive presentation satisfying
$F(1)=1$, $F(x+1)=xF(x)$, and log-convexity, with local positive bounds so
that $f=\log F$ is defined.  Choose an integer $B\ge\max(1,b)$ and an
integer $n\ge2$.  Convexity, first using the points $n-1,n,n+x$, and then
the points $n,n+x,n+B$, gives
\[
 f(n)+x\log(n-1)
 \le f(n+x)
 \le f(n)+\frac{x}{B}\sum_{j=0}^{B-1}\log(n+j)
 \le f(n)+x\log(n+B-1).
\]
Here $f(n)=\log((n-1)!)$, and the recurrence gives
$f(n+x)=f(x)+\sum_{k=0}^{n-1}\log(x+k)$.  Therefore $f(x)$ belongs to
the explicit interval
\begin{align*}
 L_n(x)&=\log((n-1)!)+x\log(n-1)
              -\sum_{k=0}^{n-1}\log(x+k),\\
 U_n(x)&=\log((n-1)!)+x\log(n+B-1)
              -\sum_{k=0}^{n-1}\log(x+k).
\end{align*}
Its width satisfies
\begin{equation}\label{eq:bm-width-general}
 0\le U_n(x)-L_n(x)
 =x\log\left(1+\frac{B}{n-1}\right)
 \le\frac{B^2}{n-1}.
\end{equation}
The already constructed gamma presentation satisfies the same hypotheses,
so both $f(x)$ and $g(x)$ lie in this interval.  Choosing
$n=1+B^2 2^{p+2}$ proves $|f(x)-g(x)|\le2^{-p-2}$, apart from elementary
function approximation errors which are assigned the remaining budget.
Thus $F(x)=\Gamma(x)$ as equality of names.  Exponentiation is legitimate
with the explicit local bound supplied either in the presentation of $F$
or by the displayed squeeze at $n=2$.

This proof uses only two convexity instances and $n+B$ recurrence instances
at each requested precision.  It does not find the integer part of an
arbitrary real, nor use a completeness principle for arbitrary functions.
For rational arguments, log-convexity may equivalently be supplied by the
integer-power inequality
$F((u x+v y)/(u+v))^{u+v}\le F(x)^uF(y)^v$.

\subsection{Wendel, Gautschi, and effective gamma-ratio asymptotics}

For $x>0$ and $0\le s\le1$, log-convexity between $x$ and $x+1$ gives
$\Gamma(x+s)\le x^s\Gamma(x)$.  Applying the same inequality to $x+s$
with increment $1-s$ gives
$x\Gamma(x)\le(x+s)^{1-s}\Gamma(x+s)$.  Thus
\begin{equation}\label{eq:wendel-constructive}
 \left(\frac{x}{x+s}\right)^{1-s}
 \le\frac{\Gamma(x+s)}{x^s\Gamma(x)}\le1.
\end{equation}
For $0<s<1$ both inequalities are strict, with gaps obtainable from
\eqref{eq:strong-log-convexity}.  This also gives Gautschi's inequalities
\begin{equation}\label{eq:gautschi-constructive}
 x^{1-s}<\frac{\Gamma(x+1)}{\Gamma(x+s)}<(x+1)^{1-s}
 \qquad(0<s<1).
\end{equation}
Indeed the first follows from the preceding upper bound for
$\Gamma(x+s)$; the preceding second application gives the stronger upper
bound $(x+s)^{1-s}<(x+1)^{1-s}$.

Writing the middle expression of \eqref{eq:wendel-constructive} as $Q$,
\[
 -\frac{s(1-s)}x\le\log Q\le0,
 \qquad 0\le1-Q\le\frac{s(1-s)}x\le\frac1{4x}.
\]
This is an explicit uniform gamma-ratio asymptotic, not just a limit
assertion.  More generally, let $|\alpha|,|\beta|\le A$ with $A\ge1$ and
$x\ge2A+1$.  On the segment joining $x+\alpha$ to $x+\beta$,
\eqref{eq:digamma-elementary-bounds} implies
\[
 |\psi(x+t)-\log x|
 \le\frac{2(A+1)}x\quad(|t|\le A).
\]
Integration over this finite segment gives
\begin{equation}\label{eq:general-gamma-ratio}
 \left|\log\frac{\Gamma(x+\alpha)}{\Gamma(x+\beta)x^{\alpha-\beta}}\right|
 \le\frac{4A(A+1)}x.
\end{equation}
If $C=4A(A+1)$ and $x\ge\max(2A+1,C)$, the relative error is at most
$2C/x$, since $|e^u-1|\le2|u|$ for $|u|\le1$.

\section{Critical points, including every negative interval}
\label{sec:gamma-critical-points}

The standard harmonic approximation to Euler's constant satisfies
\[
 H_N-\log(N+1)\le\gamma\le H_N-\log N,
\]
with bracket width at most $1/N$.  The elementary series
$\log2=2\sum_{j\ge0}3^{-2j-1}/(2j+1)$ gives
$2/3<\log2<3/4$.  Taking $N=1$ in the lower bracket and $N=4$ in the
upper bracket proves
\[
 \frac14<\gamma<\frac34.
\]
The digamma series therefore gives
$\psi(1)=-\gamma<-1/4$ and $\psi(2)=1-\gamma>1/4$.
Moreover $1/4\le\psi'(x)\le2$ on $[1,2]$.

Here is an explicit construction of the zero, avoiding a decision whether
an arbitrary computable real is positive.  Maintain rational endpoints
$l,u$ with $\psi(l)\le0\le\psi(u)$ and let $w=u-l$, $m=(l+u)/2$.
Compute a rational $q$ with $|q-\psi(m)|\le\varepsilon=w/64$.
If $q>\varepsilon$, replace $u$ by $m$; if $q<-\varepsilon$, replace $l$
by $m$.  In the remaining, decidable rational case, $|\psi(m)|\le2\varepsilon$.
The derivative lower bound shows that the interval
\[
 [l,u]\cap[m-8\varepsilon,m+8\varepsilon]
\]
still has the required endpoint signs.  Replace $[l,u]$ by that interval.
The width decreases by at least a factor of two each time.  All tests
are rational, all approximations have predetermined error budgets, and
iteration is primitive recursion.  The nested midpoints consequently give
a primitive recursive name $\alpha$ with $1<\alpha<2$ and $\psi(\alpha)=0$.
For the last equality, the bound $\psi'\le2$ bounds the value error by
twice the interval width.  Strict monotonicity of $\psi$ proves uniqueness.
Since $\Gamma'=\Gamma\psi$ and gamma is positive, this point is its unique
global minimum on the positive half-line.  To make this last conclusion
without real-order trichotomy, use the integrated tangent inequality
\[
 g(y)-g(v)-g'(v)(y-v)
 =(y-v)^2\int_0^1(1-t)g''(v+t(y-v))\,dt
 \ge\frac{(y-v)^2}{2b^2}
\]
on any positive bounded interval containing $y$ and $v$.  Substitute
$v=\alpha$ and $g'(\alpha)=0$.  The explicit square gives strictness when
$|y-\alpha|$ has a positive rational lower bound.

The same construction works on every interval $(-n-1,-n)$, $n\ge0$.
For $x$ in that interval, recurrence gives
\begin{equation}\label{eq:negative-digamma}
 \psi(x)=\psi(x+n+2)-\sum_{k=0}^{n+1}\frac1{x+k},\qquad
 \psi'(x)=\psi'(x+n+2)+\sum_{k=0}^{n+1}\frac1{(x+k)^2}>\frac14.
\end{equation}
Take $\delta=1/(4(n+2))$, $l=-n-1+\delta$, and $u=-n-\delta$.
At $l$ the singular term is $-1/\delta$ and at $u$ it is $+1/\delta$.
All other reciprocal terms have absolute value at most $2$; the remaining
digamma term has absolute value at most $1$.  Therefore
\[
 \psi(l)\le-1/\delta+(2n+3)<0,
 \qquad\psi(u)\ge1/\delta-(2n+3)>0.
\]
On $[l,u]$, an upper derivative bound is $2+(n+2)\delta^{-2}$.
The preceding rational algorithm constructs the unique zero $\alpha_n$.
The sign of gamma on the whole interval is $(-1)^{n+1}$, by its recurrence
to the positive interval $(1,2)$.  Thus $\alpha_n$ is a local maximum on
each negative-sign interval and a local minimum on each positive-sign
interval.  In fact these are extrema on the entire indicated intervals:
apply the same tangent inequality to $\log|\Gamma|$, whose second
derivative is $\psi'>0$, and then use the fixed sign of gamma.  This
again avoids splitting an arbitrary real according to its order relative
to the critical point.  No other real critical points exist away from
the poles.

\section{Stirling's formula from finite differences}
\label{sec:stirling-differences}

Define
\[
 A(x)=\log\Gamma(x)-(x-\tfrac12)\log x+x,\qquad
 C=\tfrac12\log(2\pi).
\]
The recurrence for gamma gives the exact finite-difference identity
\[
 d(x):=A(x)-A(x+1)
 =(x+\tfrac12)\log(1+1/x)-1.
\]
With $z=(2x+1)^{-1}\in(0,1)$, the finite geometric expansion for the
logarithm gives
\begin{equation}\label{eq:stirling-d-series}
 d(x)=\sum_{j=1}^{\infty}\frac{z^{2j}}{2j+1}.
\end{equation}
After $j=K$, the error is at most
$z^{2K+2}/((2K+3)(1-z^2))$.  In particular, retaining the improvement in
the coefficient at $j=2$,
\begin{equation}\label{eq:stirling-d-telescoper}
 \frac1{3(2x+1)^2}\le d(x)
 \le\frac1{12x(x+1)}-\frac2{15(2x+1)^4}.
\end{equation}
The upper main term is the difference $1/(12x)-1/(12(x+1))$.
Hence for integers $M>N\ge1$,
\[
 0\le A(N)-A(M)\le\frac1{12N}-\frac1{12M}.
\]
This is a finite sum of inequalities.  It constructs, for example by
sampling $N=2^{p+4}$, a regular primitive recursive name $c$ with
$|A(N)-c|\le1/(12N)$.  No arbitrary monotone-convergence theorem is used.

From \eqref{eq:digamma-elementary-bounds},
\[
 -\frac1{2x}\le A'(x)=\psi(x)-\log x+\frac1{2x}\le0.
\]
Thus for $0\le v\le b$, $|A(N+v)-A(N)|\le b/(2N)$.  In particular the
same $c$ is approached along every fixed translated sequence $N+x$,
with an explicit rational rate.

The duplication formula, proved independently from the finite gamma
products and reflection, gives
\begin{equation}\label{eq:constant-by-duplication}
 A(x)+A(x+\tfrac12)-A(2x)
 =C+\tfrac12-x\log(1+1/(2x)).
\end{equation}
The last two terms lie between $0$ and $1/(8x)$, by
$u-u^2/2\le\log(1+u)\le u$.  Setting $x=N$ and using the preceding
Cauchy estimates proves $|c-C|\le7/(12N)$ for every $N$.  Substituting a
sufficiently large explicit power of two for $N$ proves equality of the
two names.  This determines the constant without choosing a subsequence
or an unknown limiting value.

We now have
\begin{equation}\label{eq:stirling-positive-remainder}
 R(x):=A(x)-C=\sum_{k=0}^{\infty}d(x+k),\qquad
 0\le\sum_{k=N}^{\infty}d(x+k)\le\frac1{12(x+N)}.
\end{equation}
The equality means that for a requested precision the two finite
expressions differ by the displayed tail.  Keeping the first positive
margin in \eqref{eq:stirling-d-telescoper} proves
\begin{equation}\label{eq:stirling-first-bound}
 \frac1{3(2x+1)^2}\le R(x)
 \le\frac1{12x}-\frac2{15(2x+1)^4}<\frac1{12x}
 \qquad(x>0).
\end{equation}
In particular
\begin{equation}\label{eq:stirling-multiplicative-bound}
 1<\frac{\Gamma(x)}{\sqrt{2\pi}\,x^{x-1/2}e^{-x}}<e^{1/(12x)}.
\end{equation}
For $x\ge1$, the relative error is at most $1/(6x)$, since
$e^u-1\le2u$ when $0\le u\le1$.

The sharper Robbins lower bound also follows by rational telescoping.
For $x\ge1$ put
\[
 q(x)=\frac1{12x+1},\qquad
 \lambda(x)=\frac{24x-23}
 {3(2x+1)^2(12x+1)(12x+13)}>0.
\]
The polynomial identity
$(12x+1)(12x+13)-36(2x+1)^2=24x-23$ gives
$d(x)\ge q(x)-q(x+1)+\lambda(x)$.  Summing finitely and using
\eqref{eq:stirling-positive-remainder} proves
\begin{equation}\label{eq:robbins-real}
 \frac1{12x+1}+\lambda(x)\le R(x)
 \le\frac1{12x}-\frac2{15(2x+1)^4}\qquad(x\ge1).
\end{equation}
Taking $x=n$ and using $\Gamma(n+1)=n!$ gives the usual strict factorial
Robbins inequalities.  This argument is the real-variable extension of
the checked finite telescoper in the repository's Stirling draft.

\section{All orders of Stirling's expansion with the sharp signed bound}
\label{sec:all-order-stirling}

The order $m\ge1$ may be fixed externally; the following is then a family
of arithmetic theorems with primitive recursive precision parameters.
Bernoulli polynomials have rational coefficients and are defined without
any analytic generating function: $B_0(t)=1$, $B_r'=rB_{r-1}$, and
$\int_0^1 B_r(t)\,dt=0$ for $r\ge1$.  Antidifferentiating a polynomial
and subtracting its mean is a finite rational algorithm.  Write
$b_r=B_r(0)$.  Uniqueness of these elementary polynomial conditions gives
\[
 B_r(1-t)=(-1)^r B_r(t),\qquad
 B_r(1)=B_r(0)\quad(r\ge2),\qquad b_{2j+1}=0\quad(j\ge1).
\]
For example $B_1=t-1/2$, $B_2=t^2-t+1/6$,
$B_4=t^4-2t^3+t^2-1/30$, and $b_6=1/42$.

\subsection{Finite Euler--Maclaurin, derived rather than assumed}

On $[k,k+1]$ set $\widetilde B_r(t)=B_r(t-k)$, using the common endpoint
value for even $r$.  Integration by parts on that interval gives
\[
 \int_k^{k+1}\widetilde B_1(t)f'(t)\,dt
 =\frac{f(k)+f(k+1)}2-\int_k^{k+1}f(t)\,dt.
\]
Sum over $k=0,\ldots,N-1$ and integrate the last term successively using
$B_r'=rB_{r-1}$.  Odd endpoint contributions vanish and even endpoint
contributions telescope.  One obtains the finite identity
\begin{align}
 \sum_{k=0}^{N}f(k)
 ={}&\int_0^N f(t)\,dt+\frac{f(0)+f(N)}2\notag\\
 &+\sum_{r=1}^{m}\frac{b_{2r}}{(2r)!}
       \bigl(f^{(2r-1)}(N)-f^{(2r-1)}(0)\bigr)
 -\frac1{(2m)!}\int_0^N\widetilde B_{2m}(t)f^{(2m)}(t)\,dt.
 \label{eq:finite-em-gamma}
\end{align}
For the functions used here, integration by parts itself follows by
finite telescoping product increments; derivatives and continuity bounds
are explicit rational functions on a compact interval.  Thus this
identity introduces no nonconstructive calculus principle.

Apply \eqref{eq:finite-em-gamma} to $f(t)=(x+t)^{-1}$ and use the finite
Euler expression for $\psi$.  Since
$f^{(j)}(t)=(-1)^j j!(x+t)^{-j-1}$, all terminal errors vanish with
explicit bounds.  The resulting identity is
\begin{equation}\label{eq:psi-em-exact}
 \psi(x)=\log x-\frac1{2x}
 -\sum_{r=1}^{m}\frac{b_{2r}}{2r x^{2r}}
 +\int_0^\infty\frac{\widetilde B_{2m}(t)}{(x+t)^{2m+1}}\,dt.
\end{equation}
If $M_m$ is the sum of the absolute values of the coefficients of
$B_{2m}$, the last integral after a cutoff $T$ has tail at most
$M_m/(2m(x+T)^{2m})$.  All terminal logarithms are bounded using
$\log(1+x/N)\le b/N$.

Integrate \eqref{eq:psi-em-exact} from $x$ to a finite $Y$, use
$A(Y)-C\to0$ from the preceding explicit Stirling bound, and regroup the
finite double sums.  Integrating the rational kernel gives
\begin{equation}\label{eq:loggamma-em-exact}
 R(x)=\sum_{r=1}^{m}\frac{b_{2r}}{2r(2r-1)x^{2r-1}}
 -\frac1{2m}\int_0^\infty
             \frac{\widetilde B_{2m}(t)}{(x+t)^{2m}}\,dt.
\end{equation}
The omitted rectangular tails in this double-sum argument are at most
$M_m/[2m(2m-1)(x+T)^{2m-1}]$ and
$M_m/[2m(2m-1)Y^{2m-1}]$, respectively.  Thus a general Fubini theorem
is unnecessary.

\subsection{The Bernoulli sign and a finite reverse-Chebyshev argument}

We prove the elementary shape fact needed for the sharp remainder.
On $0<t<1/2$, $B_1(t)<0$.  For $j\ge1$ the polynomial $B_{2j+1}$ vanishes
at $0$ and $1/2$, and
$B_{2j+1}''=(2j+1)(2j)B_{2j-1}$.  Induction therefore proves
\[
 (-1)^jB_{2j-1}(t)>0\qquad(0<t<1/2).
\]
Indeed a function with positive second derivative and zero endpoints is
negative between them, and with negative second derivative it is
positive; for these polynomials this statement follows from the finite
secant inequality proved by integrating derivative increments.

It follows that
\[
 H_m(t)=(-1)^m(B_{2m}(t)-b_{2m})
\]
is symmetric about $1/2$, is zero at $0$ and $1$, is positive in between,
and is strictly increasing on $[0,1/2]$.  Its mean is
\[
 \int_0^1H_m(t)\,dt=(-1)^{m+1}b_{2m}=|b_{2m}|>0.
\]
This also proves the alternating signs of the nonzero even Bernoulli
numbers.

Fix $A>0$ and $q>0$, and put $w(t)=(A+t)^{-q}$.  The paired weight
$W(t)=w(t)+w(1-t)$ is decreasing for $0\le t\le1/2$, since
\[
 W'(t)=-q(A+t)^{-q-1}+q(A+1-t)^{-q-1}\le0.
\]
For an increasing function $H$ and decreasing function $W$ on an interval
of length $L$, the finite double-sum identity for
$(H(u)-H(v))(W(u)-W(v))\le0$ gives
\[
 \int HW\le L^{-1}\left(\int H\right)\left(\int W\right).
\]
Passing from the finite sums uses only their uniform moduli.  Apply this
on $[0,1/2]$ to $H_m$ and the paired weight.  Symmetry yields
\begin{equation}\label{eq:bernoulli-weight-bound}
 0<\int_0^1H_m(t)(A+t)^{-q}\,dt
 <|b_{2m}|\int_0^1(A+t)^{-q}\,dt.
\end{equation}
Strictness has constructive content: $H_m$ is strictly increasing and
$W$ strictly decreasing on the interior, so two disjoint rational
subintervals give an explicit negative margin in the double-sum
identity.  Their rational centers can be fixed, for example at $1/8$
and $3/8$, and continuity bounds shrink the intervals to preserve the
strict gap.  The coefficients of $H_m$ and bounds for $A$ compute the
required widths; no unbounded search is involved.
For an explicit weight gap, if $A\le B$ and $t\in[1/8,3/8]$, then
\[
 -W'(t)\ge\frac{q(q+1)}{4(B+1)^{q+2}},\qquad
 W(1/8)-W(3/8)\ge\frac{q(q+1)}{16(B+1)^{q+2}}.
\]
The other gap, $H_m(3/8)-H_m(1/8)>0$, is an exactly computed rational.

Subtract the first $m-1$ correction terms in
\eqref{eq:loggamma-em-exact}.  Since
$\int_0^\infty(x+t)^{-2m}\,dt=x^{1-2m}/(2m-1)$, the result is
\begin{equation}\label{eq:stirling-remainder-positive-integral}
 E_m(x):=R(x)-\sum_{r=1}^{m-1}
            \frac{b_{2r}}{2r(2r-1)x^{2r-1}}
 =\frac{(-1)^{m+1}}{2m}
       \int_0^\infty\frac{H_m(\{t\})}{(x+t)^{2m}}\,dt.
\end{equation}
Here $H_m(\{t\})$ is implemented piecewise on explicitly indexed unit
intervals; computing the floor of an arbitrary real is not required.
Apply \eqref{eq:bernoulli-weight-bound} on each unit interval and add the
finite inequalities.  The integral tail is explicit.  Keeping the
positive gap from the first interval proves the sharp signed bound
\begin{equation}\label{eq:sharp-stirling-remainder}
 \boxed{\quad
 0<(-1)^{m+1}E_m(x)
 <\frac{|b_{2m}|}{2m(2m-1)x^{2m-1}}
 \quad(x>0).\quad}
\end{equation}
Consequently
\[
 \log\Gamma(x)\sim
 (x-\tfrac12)\log x-x+\tfrac12\log(2\pi)
 +\frac1{12x}-\frac1{360x^3}+\frac1{1260x^5}-\cdots,
\]
where the meaning of every truncation is precisely
\eqref{eq:sharp-stirling-remainder}.  For example,
\[
 \frac1{12x}-\frac1{360x^3}<R(x)<\frac1{12x},\qquad
 0<R(x)-\frac1{12x}+\frac1{360x^3}<\frac1{1260x^5}.
\]

The same argument applied to \eqref{eq:psi-em-exact} gives, if
\[
 D_m(x)=\psi(x)-\log x+\frac1{2x}
             +\sum_{r=1}^{m-1}\frac{b_{2r}}{2r x^{2r}},
\]
the derivative expansion with the sharp signed bound
\begin{equation}\label{eq:sharp-digamma-remainder}
 0<(-1)^mD_m(x)<\frac{|b_{2m}|}{2m x^{2m}}.
\end{equation}

Exponentiating finite truncations gives the multiplicative Stirling
expansion.  Its rational coefficients are determined primitively by
$c_0=1$ and
\[
 n c_n=\sum_{j=1}^{n}j\ell_jc_{n-j},\qquad
 \ell_{2r-1}=\frac{b_{2r}}{2r(2r-1)},\quad \ell_{2r}=0.
\]
Thus the first terms are
\[
 \frac{\Gamma(x)}{\sqrt{2\pi}\,x^{x-1/2}e^{-x}}
 \sim1+\frac1{12x}+\frac1{288x^2}
          -\frac{139}{51840x^3}-\frac{571}{2488320x^4}+\cdots.
\]
For any requested fixed truncation order, choose a log-gamma truncation
with error of higher degree, substitute its finite polynomial into the
finite exponential sum, and bound the two remainders with
$|e^u-e^v|\le e^{\max(|u|,|v|)}|u-v|$ and the geometric majorant for the
exponential Taylor tail.  This is an explicit finite procedure for the
error constant, not an appeal to closure properties of an undefined
asymptotic notation.

\subsection{Primitive recursive quadrature budgets}

Here are sufficient budgets for every improper integral introduced in
this section.  If $P(t)$ is one of the fixed periodic Bernoulli
polynomials, let $M$ be the sum of its absolute coefficient values and
$L$ the sum of the absolute coefficient values of its derivative.  For
$q\ge2$ and $x\ge a>0$, the function
$P(\{t\})(x+t)^{-q}$ is globally Lipschitz on $[0,T]$ with constant
\[
 K=L a^{-q}+qM a^{-q-1}.
\]
The pieces join continuously because the endpoint values of the even
Bernoulli polynomials agree.  A left sum with $J$ equal subintervals has
error at most $KT^2/(2J)$.  Its tail after $T\ge1$ is at most
$M/((q-1)(x+T)^{q-1})\le M/T$.
For target error $2^{-p}$ one may therefore take
\[
 T=\max(1,\lceil 2^{p+3}M\rceil),\qquad
 J=\max(1,\lceil 2^{p+3}KT^2\rceil),
\]
and allocate the remaining error to approximating $x$ in the summands.
The floor symbols in the notation are computed only for rational mesh
points, by integer division.  All loops are bounded, all coefficients
are rational, and all required sums and products are primitive recursive.
For each fixed derivative or asymptotic order, these give the explicit
finite rational identities and inequalities to be formalized inside iPRA.

% Sources for root's bibliography:
% NIST DLMF §§5.5(iv), 5.6(i), 5.11(i)-(iii).
% Repository commit 94a4713736d281e1bfcd0b54e0acb6410136971c:
% bohr_mollerup blob 55a972a2cd46f9b4a5edd70b79598cca3413ccd7;
% stirling blob 6c6f8e17747d567b426b19f0cb2678cba303457f.
